\documentclass[11pt]{article}
\usepackage[margin=1in]{geometry}
\usepackage{amsmath,amssymb,amsthm,mathtools}
\usepackage{graphicx}
\usepackage[hidelinks]{hyperref}
\usepackage{microtype}

\newtheorem{theorem}{Theorem}
\newtheorem{proposition}[theorem]{Proposition}
\newtheorem{remark}[theorem]{Remark}
\newcommand{\Hh}{\mathcal H}
\newcommand{\N}{\mathbb N}

\title{The Unilateral Shift Obstructs the Nontrivial\\
Frame-Orbit Decomposition Conjecture\\
\large A scoped counterexample to the final conjecture of arXiv:2301.06183}
\author{Automated mathematics run \texttt{fa\_banach\_001}, lane 9}
\date{27 August 2026}

\begin{document}
\maketitle

\begin{center}
\fbox{\parbox{0.91\textwidth}{\textbf{Status.}
Candidate counterexample, likely valid, pending expert review.  The printed
statement is tautological if zero summands are allowed; its natural nontrivial
reading, requiring at least two nonzero summands, is false.}}
\end{center}

\section{Source conjecture}

For a separable Hilbert space $\Hh$, the source paper
\cite{CheshmavarDallakiBaradaran} defines
\[
 E(\Hh)=\left\{T\in B(\Hh):
 \{T^k\varphi\}_{k=0}^{\infty}\text{ is a frame for $\Hh$ for
 some $\varphi\in\Hh$}\right\}.
\]
On page 11, its final conjecture states that, for every $T\in E(\Hh)$,
there is a Hilbert direct-sum decomposition
\[
 \Hh=\bigoplus_{n\in\N}\Hh_n
\]
such that (i) $T\Hh_n\subseteq\Hh_n$, and (ii) for every $n$ there
is $\varphi_n\in\Hh_n$ for which
$\{T^k\varphi_n\}_{k=0}^{\infty}$ is a frame for $\Hh_n$.

\begin{figure}[ht]
  \centering
  \includegraphics[width=\textwidth]{figures/open_problem_crop.png}
  \caption{The final conjecture on page 11 of the source paper.}
\end{figure}

\section{Idea of the counterexample}

The unilateral shift is already the canonical example of an operator with a
frame orbit: the orbit of its first basis vector is the entire standard
orthonormal basis.  At the same time, the multiplicity-one shift cannot split
into two nonzero invariant blocks.  Indeed, in an orthogonal decomposition
into invariant summands, the complement of each summand is invariant as well,
so its orthogonal projection commutes with the shift.  But the only orthogonal
projections commuting with the multiplicity-one unilateral shift are $0$ and
the identity.

This proves that the only possible direct-sum presentation has one nonzero
summand, namely the whole space.  Thus the conjecture has a sharp convention
dichotomy: it is automatic if zero summands are admitted and false if the
summands are required to give a genuine decomposition.

\section{The literal convention}

\begin{proposition}[Vacuity of the literal statement]\label{prop:literal}
Suppose zero Hilbert spaces are allowed as direct summands and the zero orbit
is regarded as a frame for the zero space.  Then the printed conjecture is
true immediately for every $T\in E(\Hh)$.
\end{proposition}

\begin{proof}
By the definition of $E(\Hh)$ there exists $\varphi\in\Hh$ such that
$\{T^k\varphi\}_{k\ge0}$ is a frame for $\Hh$.  Set
\[
 \Hh_1=\Hh,\qquad \varphi_1=\varphi,
 \qquad \Hh_n=\{0\},\quad \varphi_n=0\quad(n\ge2).
\]
Then $\Hh=\bigoplus_{n\in\N}\Hh_n$, every summand is invariant,
and the frame inequalities on a zero space are vacuous.
\end{proof}

If the convention is instead that every displayed summand must be nonzero,
the next theorem supplies a counterexample.

\section{The unilateral-shift counterexample}

Let $\Hh=\ell^2(\N_0)$ with standard orthonormal basis
$(e_j)_{j\ge0}$, and let $S$ be the unilateral shift
\[
 Se_j=e_{j+1}\qquad(j\ge0).
\]

\begin{theorem}[Failure of every nontrivial decomposition]
\label{thm:counterexample}
The operator $S$ belongs to $E(\Hh)$, but $\Hh$ cannot be written as an
orthogonal direct sum of two or more nonzero closed $S$-invariant subspaces.
Consequently, the source conjecture is false under the natural interpretation
that its countable family of summands is nonzero.
\end{theorem}

\begin{proof}
First,
\[
 \{S^ke_0\}_{k=0}^{\infty}=\{e_k\}_{k=0}^{\infty}
\]
is an orthonormal basis for $\Hh$, hence a Parseval frame.  Therefore
$S\in E(\Hh)$.

Suppose that
\[
 \Hh=\bigoplus_{n\in I}\Hh_n
\]
is an orthogonal direct sum of closed $S$-invariant subspaces.  Fix $n$ and
let $P_n$ be the orthogonal projection onto $\Hh_n$.  The orthogonal
complement
\[
 \Hh_n^\perp=\bigoplus_{m\ne n}\Hh_m
\]
is also $S$-invariant: finite partial sums are carried into the same
complement, and boundedness of $S$ passes this to the closure.  Thus $S$ is
block diagonal relative to $\Hh_n\oplus\Hh_n^\perp$, so
\begin{equation}\label{eq:commute}
 P_nS=SP_n.
\end{equation}
Because $P_n=P_n^*$, taking adjoints in \eqref{eq:commute} also gives
$S^*P_n=P_nS^*$.  Now
\[
 \ker S^*=\operatorname{span}\{e_0\},
\]
so $P_ne_0=ce_0$ for some scalar $c$.  Since $P_n$ is a projection,
$c^2=c$, hence $c\in\{0,1\}$.  Finally, for every $k\ge0$,
\[
 P_ne_k=P_nS^ke_0=S^kP_ne_0=ce_k.
\]
Therefore $P_n=cI$.  Each $\Hh_n$ is consequently either $\{0\}$ or
$\Hh$, and there cannot be two nonzero summands.  In particular, no
countably infinite nonzero decomposition satisfying condition (i) exists;
condition (ii) cannot repair this obstruction.
\end{proof}

\section{Even nonorthogonal topological sums fail}

The use of ``direct sum of Hilbert spaces'' normally denotes an orthogonal
Hilbert direct sum, which Theorem \ref{thm:counterexample} handles.  The same
counterexample survives a broader bounded topological-direct-sum reading.

\begin{proposition}\label{prop:topological}
The commutant of the multiplicity-one unilateral shift contains no
nontrivial idempotent.  Hence $\ell^2(\N_0)$ cannot be a bounded topological
direct sum of two nonzero closed $S$-invariant subspaces either.
\end{proposition}

\begin{proof}
Identify $\ell^2(\N_0)$ unitarily with the Hardy space $H^2(\mathbb D)$,
so that $S$ becomes multiplication $M_z$ by the coordinate function.  If a
bounded operator $A$ commutes with $M_z$, let $\psi=A1$.  For every polynomial
$q$,
\[
 Aq=Aq(M_z)1=q(M_z)A1=q\psi.
\]
Boundedness and density show that $A=M_\psi$ for a bounded analytic
multiplier $\psi\in H^\infty(\mathbb D)$.  If $A^2=A$, then
$\psi^2=\psi$.  The analytic function $\psi$ takes values only in
$\{0,1\}$, so connectedness of $\mathbb D$ forces it to be constant.  Thus
$A=0$ or $A=I$.

If $\Hh=M\dotplus N$ were a bounded topological direct sum of two closed
$S$-invariant subspaces, the coordinate projection onto $M$ along $N$ would
be a bounded idempotent commuting with $S$, contradicting the preceding
paragraph unless one summand were zero.
\end{proof}

\section{Verification and limitations}

The counterexample is scoped to the only nontrivial reading of the printed
conjecture: at least two nonzero summands.  It does not claim that the literal
zero-summand convention is false; Proposition \ref{prop:literal} shows that
version is automatic.

The script \texttt{code/verify\_shift\_obstruction.py} checks the analogous
matrix statements for truncated shifts in every dimension from 2 through 40.
It verifies that the first-vector orbit has identity Gram matrix and that the
joint commutant of the shift and its adjoint is one-dimensional.  These 39
finite checks are designed to catch orientation and commutator mistakes; the
infinite-dimensional proof above is exact and does not depend on them.

The most important human-review question is semantic rather than analytic:
whether the source authors intended all $\Hh_n$ to be nonzero.  Under that
intention, Theorem \ref{thm:counterexample} is a complete disproof.  Without
it, Proposition \ref{prop:literal} is a complete literal answer.

\section{Bounded novelty check}

On 27 August 2026, the run registry, solution, attempt, and proof-gap indexes
were searched for arXiv:2301.06183, its title, unilateral-shift frame orbits,
invariant-summand decompositions, and reducing-subspace variants.  Bounded
exact-phrase, arXiv, title, and citation searches found the source paper, the
related arXiv papers 2212.01921 and 2505.19303, and a 2025 article citing the
source (doi:10.1007/s11785-025-01839-8).  No exact answer to the final
conjecture surfaced; the available metadata and abstract of the citing paper
concern other characterizations.  Novelty is plausible, not certified.

\begin{thebibliography}{9}
\bibitem{CheshmavarDallakiBaradaran}
J. Cheshmavar, A. Dallaki, and J. Baradaran,
\emph{On abstract results of operator representation of frames in Hilbert spaces},
J. Pseudo-Differ. Oper. Appl. 14 (2023), article 30,
arXiv:2301.06183v1, doi:10.1007/s11868-023-00524-8.
\end{thebibliography}

\end{document}
